\documentclass[a4paper,12pt]{article}

\usepackage[T2A]{fontenc}
\usepackage[utf8]{inputenc}
\usepackage[russian]{babel}

\usepackage{amsmath,amssymb,amsthm,mathtools}
\usepackage{helvet}
\renewcommand{\familydefault}{\sfdefault}

\usepackage{icomma}
\usepackage[a4paper,margin=20mm,headheight=15pt]{geometry}
\usepackage{microtype}

\usepackage{tikz}
\usetikzlibrary{calc}
\usetikzlibrary{arrows.meta,positioning,shapes.geometric}

\usepackage{tkz-euclide}

\usepackage{pgfplots}
\pgfplotsset{compat=1.18}

\usepackage{tabularx,booktabs,longtable}
\usepackage{enumitem}
\usepackage{hyperref}
\usepackage{parskip}
\usepackage{fancyhdr}
\usepackage{setspace}
\usepackage{xcolor}

\definecolor{accent}{HTML}{008080}
\definecolor{darkaccent}{HTML}{004D4D}
\definecolor{textgray}{HTML}{333333}
\definecolor{softgray}{HTML}{F1F4F4}

\color{textgray}

\onehalfspacing
\setlength{\parindent}{0pt}
\setlength{\parskip}{0.55em}

\hypersetup{
    colorlinks=true,
    linkcolor=darkaccent,
    urlcolor=darkaccent
}

\newcommand{\tg}{\mathop{\rm tg}\nolimits}
\newcommand{\ctg}{\mathop{\rm ctg}\nolimits}
\newcommand{\arctg}{\mathop{\rm arctg}\nolimits}

\pagestyle{fancy}
\fancyhf{}
\fancyhead[L]{\small Чек-лист: графики функций}
\fancyhead[R]{\small Матвей Горбачёв}
\fancyfoot[C]{\thepage}
\renewcommand{\headrulewidth}{0.4pt}

\newcommand{\sectionline}{%
    \vspace{0.2em}
    {\color{accent}\rule{\linewidth}{0.7pt}}
    \vspace{0.45em}
}

\newcommand{\key}[1]{\textcolor{darkaccent}{\textbf{#1}}}

\tikzset{
    flowstart/.style={
        ellipse,
        draw=darkaccent,
        line width=0.8pt,
        minimum width=42mm,
        minimum height=11mm,
        align=center,
        font=\small
    },
    flowaction/.style={
        rectangle,
        rounded corners=2mm,
        draw=accent,
        line width=0.8pt,
        minimum width=78mm,
        text width=70mm,
        minimum height=13mm,
        align=center,
        font=\small
    },
    flowdecision/.style={
        diamond,
        aspect=2.35,
        draw=darkaccent,
        line width=0.8pt,
        text width=51mm,
        align=center,
        inner sep=1.6mm,
        font=\small
    },
    flowarrow/.style={
        -{Stealth[length=3mm,width=2mm]},
        line width=0.8pt,
        draw=darkaccent
    }
}

\begin{document}

% =========================================================
% ТИТУЛЬНЫЙ ЛИСТ
% =========================================================

\begin{titlepage}
\thispagestyle{empty}

\begin{center}

\vspace*{22mm}

{\color{darkaccent}\Large\bfseries ЧЕК-ЛИСТ}

\vspace{5mm}

{\Huge\bfseries
Графики функций\\[2mm]
и метод узловых точек}

\vspace{9mm}

{\Large
Поиск коэффициентов по графику\\
и вычисление значений функции}

\vspace{19mm}

{\large
Матвей Горбачёв}

\vspace{5mm}

{\normalsize
Подготовка к ЕГЭ по профильной математике}

\vfill

{\large
\textbf{Лёвин Артём Александрович}\\[2mm]
эксклюзивно для Матвея Горбачёва}

\vspace{7mm}

{\large \today}

\vspace{18mm}

\begin{tikzpicture}[remember picture,overlay]
    \draw[
        accent,
        line width=1.1pt
    ]
    ($(current page.south west)+(18mm,16mm)$)
    .. controls
    ($(current page.south west)+(65mm,25mm)$)
    and
    ($(current page.south east)+(-65mm,7mm)$)
    ..
    ($(current page.south east)+(-18mm,16mm)$);
\end{tikzpicture}

\end{center}
\end{titlepage}

% =========================================================
% ОСНОВНОЙ ЧЕК-ЛИСТ
% =========================================================

\section*{\textcolor{darkaccent}{1. Что требуется от задачи}}

\sectionline

В задачах по графикам часто известен \key{вид функции}, но неизвестны её коэффициенты.

Типичные формулы:

\[
f(x)=a^x+b,
\qquad
f(x)=a^{x+b},
\qquad
f(x)=a^x,
\qquad
f(x)=kx+b.
\]

По изображённому графику требуется:

\begin{itemize}[leftmargin=7mm,itemsep=2mm]
    \item восстановить неизвестные коэффициенты;
    \item записать формулу функции;
    \item вычислить значение вида \(f(c)\);
    \item найти \(x\), если задано значение \(f(x)\).
\end{itemize}

Главный универсальный инструмент --- \key{метод узловых точек}.

\section*{\textcolor{darkaccent}{2. Узловая точка}}

\sectionline

\key{Узловая точка} --- точка графика, обе координаты которой удобно считываются с координатной сетки, прежде всего целочисленные:

\[
A(x_1;y_1),\qquad B(x_2;y_2).
\]

Поскольку

\[
y=f(x),
\]

точка \(A(x_1;y_1)\), принадлежащая графику функции, даёт равенство

\[
y_1=f(x_1).
\]

Именно это позволяет превратить информацию с рисунка в уравнение.

\subsection*{\textcolor{accent}{Сколько точек нужно}}

Обычно количество выбранных узловых точек совпадает с количеством неизвестных коэффициентов.

\[
\boxed{
\text{число точек}
=
\text{число неизвестных коэффициентов}
}
\]

Например:

\[
f(x)=a^x+b
\]

содержит два неизвестных коэффициента \(a\) и \(b\), поэтому обычно достаточно двух подходящих точек.

Для

\[
f(x)=a^x
\]

неизвестен только коэффициент \(a\), поэтому достаточно одной \emph{информативной} точки.

\subsection*{\textcolor{accent}{Какие точки выбирать}}

В первую очередь удобны точки с координатами, близкими к нулю:

\[
0,\quad 1,\quad -1.
\]

Причина:

\[
a^0=1,\qquad a^1=a,
\]

поэтому система получается значительно проще.

Однако выбранная точка должна действительно содержать нужную информацию.

Например, для

\[
f(x)=a^x
\]

точка \((0;1)\) приводит к равенству

\[
1=a^0=1,
\]

из которого определить \(a\) невозможно.

В такой ситуации следует выбрать другую узловую точку.

\section*{\textcolor{darkaccent}{3. Главный алгоритм}}

\sectionline

Для функции с неизвестными коэффициентами действуйте последовательно:

\begin{enumerate}[leftmargin=9mm,itemsep=2mm]
    \item Определите вид функции.
    \item Посчитайте неизвестные коэффициенты.
    \item Выберите столько узловых точек, сколько требуется для определения коэффициентов.
    \item Подставьте координаты каждой точки в формулу:
    \[
    y=f(x).
    \]
    \item Решите полученную систему или уравнение.
    \item Проверьте ограничения на коэффициенты.
    \item Подставьте найденные коэффициенты в исходную формулу.
    \item Выполните требование задачи.
\end{enumerate}

Для показательной функции обязательно учитывайте:

\[
\boxed{a>0,\qquad a\ne1.}
\]

Это ограничение относится к основанию показательной функции.



% =========================================================
% ПРИМЕР 1
% =========================================================

\section*{\textcolor{darkaccent}{4. Пример: \(f(x)=a^x+b\)}}

\sectionline

Пусть график функции

\[
f(x)=a^x+b
\]

проходит через точки

\[
A(0;-2),
\qquad
B(1;-1).
\]

Подставляем координаты:

\[
\begin{cases}
-2=a^0+b,\\
-1=a^1+b.
\end{cases}
\]

Используем

\[
a^0=1,
\qquad
a^1=a:
\]

\[
\begin{cases}
-2=1+b,\\
-1=a+b.
\end{cases}
\]

Из первого уравнения

\[
b=-3.
\]

Тогда

\[
-1=a-3,
\]

откуда

\[
a=2.
\]

Следовательно,

\[
\boxed{f(x)=2^x-3}.
\]

Например,

\[
f(6)=2^6-3=64-3=61.
\]

\subsection*{\textcolor{accent}{Как выглядит этот график}}

\begin{center}
\begin{tikzpicture}
\begin{axis}[
    width=12.5cm,
    height=8cm,
    axis lines=middle,
    xmin=-3,
    xmax=4,
    ymin=-4,
    ymax=7,
    xtick={-3,-2,-1,0,1,2,3,4},
    ytick={-4,-3,-2,-1,0,1,2,3,4,5,6,7},
    grid=both,
    minor tick num=0,
    xlabel={$x$},
    ylabel={$y$},
    samples=250,
    clip=false,
    unit vector ratio=1 1
]
\addplot[
    domain=-3:3.25,
    very thick,
    accent
]
{2^x-3};

\addplot[
    only marks,
    mark=*,
    mark size=2.5pt,
    darkaccent
]
coordinates {(0,-2) (1,-1)};

\node[anchor=south east] at (axis cs:0,-2) {$A(0;-2)$};
\node[anchor=north west] at (axis cs:1,-1) {$B(1;-1)$};

\end{axis}
\end{tikzpicture}
\end{center}

Горизонтальная прямая

\[
y=-3
\]

является асимптотой этого графика.

% =========================================================
% ПРИМЕР 2
% =========================================================

\section*{\textcolor{darkaccent}{5. Если возникает уравнение \(a^2=2\)}}

\sectionline

Пусть после подстановки узловых точек получено

\[
a^2=2.
\]

Алгебраическое уравнение имеет два решения:

\[
a=\sqrt2,
\qquad
a=-\sqrt2.
\]

Однако для показательной функции основание должно удовлетворять условию

\[
a>0.
\]

Поэтому

\[
\boxed{a=\sqrt2}.
\]

Полезно сразу представить корень степенью:

\[
\sqrt2=2^{1/2}.
\]

Если

\[
f(x)=\left(\sqrt2\right)^x-3,
\]

то

\[
f(x)=2^{x/2}-3.
\]

Например,

\[
f(10)=2^{10/2}-3
=2^5-3
=29.
\]

\section*{\textcolor{darkaccent}{6. Степени с отрицательными показателями}}

\sectionline

Если после подстановки точки возникает

\[
a^{-2}=4,
\]

можно использовать свойство

\[
a^{-2}=\frac1{a^2}.
\]

Тогда

\[
\frac1{a^2}=4,
\qquad
a^2=\frac14.
\]

Так как \(a>0\),

\[
\boxed{a=\frac12}.
\]

Другой способ --- возвести обе части равенства в степень

\[
-\frac12:
\]

\[
a
=
4^{-1/2}
=
\left(2^2\right)^{-1/2}
=
2^{-1}
=
\frac12.
\]

Полезное правило:

\[
\boxed{(a^m)^n=a^{mn}}.
\]

% =========================================================
% f(x)=a^{x+b}
% =========================================================

\section*{\textcolor{darkaccent}{7. Если коэффициент находится внутри показателя}}

\sectionline

Рассмотрим функцию

\[
f(x)=a^{x+b}.
\]

Пусть известны точки

\[
A(-3;1),
\qquad
B(1;4).
\]

Получаем систему:

\[
\begin{cases}
1=a^{-3+b},\\
4=a^{1+b}.
\end{cases}
\]

Так как

\[
1=a^0
\]

и \(a>0,\ a\ne1\), из первого уравнения следует

\[
-3+b=0,
\]

поэтому

\[
b=3.
\]

Подставляем во второе:

\[
4=a^{1+3}=a^4.
\]

Следовательно,

\[
a=4^{1/4}.
\]

Так как

\[
4=2^2,
\]

получаем

\[
a=(2^2)^{1/4}=2^{1/2}=\sqrt2.
\]

Итак,

\[
\boxed{f(x)=\left(\sqrt2\right)^{x+3}}
\]

или

\[
\boxed{f(x)=2^{(x+3)/2}}.
\]

Например,

\[
f(-7)
=
2^{(-7+3)/2}
=
2^{-2}
=
\frac14.
\]

% =========================================================
% f(c) VS f(x)=c
% =========================================================

\section*{\textcolor{darkaccent}{8. Не путайте \(f(c)\) и \(f(x)=c\)}}

\sectionline

Это два разных типа заданий.

\subsection*{\textcolor{accent}{Тип 1. Найти \(f(3)\)}}

Если

\[
f(x)=2x+5,
\]

то в формулу вместо \(x\) подставляем \(3\):

\[
f(3)=2\cdot3+5=11.
\]

То есть

\[
\boxed{f(c)\ \Rightarrow\ x=c.}
\]

\subsection*{\textcolor{accent}{Тип 2. Найти \(x\), если \(f(x)=3\)}}

Для той же функции

\[
f(x)=2x+5
\]

записываем

\[
2x+5=3.
\]

Отсюда

\[
2x=-2,
\qquad
x=-1.
\]

То есть

\[
\boxed{f(x)=c\ \Rightarrow\ \text{значение }c
\text{ подставляется вместо }f(x).}
\]

\subsection*{\textcolor{accent}{Пример с показательной функцией}}

Если

\[
f(x)=2^x-3
\]

и требуется найти \(x\), при котором

\[
f(x)=29,
\]

то

\[
2^x-3=29,
\]

\[
2^x=32,
\]

\[
2^x=2^5.
\]

Следовательно,

\[
\boxed{x=5}.
\]

% =========================================================
% ЛИНЕЙНАЯ ФУНКЦИЯ
% =========================================================

\section*{\textcolor{darkaccent}{9. Метод работает и для линейной функции}}

\sectionline

Пусть

\[
f(x)=kx+b.
\]

Неизвестны два коэффициента:

\[
k,\qquad b.
\]

Поэтому выбираем две узловые точки.

Например, если график проходит через

\[
A(0;3),
\qquad
B(2;7),
\]

то

\[
\begin{cases}
3=k\cdot0+b,\\
7=2k+b.
\end{cases}
\]

Из первого равенства

\[
b=3.
\]

Тогда

\[
7=2k+3,
\]

\[
2k=4,
\]

\[
k=2.
\]

Следовательно,

\[
\boxed{f(x)=2x+3}.
\]

Сам алгоритм остаётся тем же:

\[
\boxed{
\text{точки}
\to
\text{подстановка}
\to
\text{система}
\to
\text{коэффициенты}
\to
\text{формула}
}
\]

% =========================================================
% ТИПИЧНЫЕ ОШИБКИ
% =========================================================

\section*{\textcolor{darkaccent}{10. Типичные ошибки}}

\sectionline

\begin{enumerate}[leftmargin=9mm,itemsep=3mm]

    \item \key{Перепутаны координаты точки.}

    Для точки

    \[
    A(2;-3)
    \]

    следует подставлять

    \[
    x=2,\qquad f(x)=-3.
    \]

    \item \key{Выбрано недостаточно точек.}

    Если неизвестны два независимых коэффициента, одного уравнения обычно недостаточно.

    \item \key{Точка не содержит информации о коэффициенте.}

    Для

    \[
    f(x)=a^x
    \]

    точка \((0;1)\) приводит только к тождеству

    \[
    1=1.
    \]

    \item \key{Забыто условие \(a>0\).}

    Если

    \[
    a^2=2,
    \]

    для основания показательной функции выбирается

    \[
    a=\sqrt2.
    \]

    \item \key{Перепутаны выражения}

    \[
    a^x+b
    \]

    и

    \[
    a^{x+b}.
    \]

    В первом случае \(b\) находится \emph{вне степени}, во втором --- \emph{в показателе степени}.

    \item \key{Неверно истолковано требование задачи.}

    \[
    f(6)
    \]

    означает подстановку \(x=6\).

    \[
    f(x)=6
    \]

    означает уравнение, которое нужно решить относительно \(x\).

    \item \key{Ошибки при работе со степенями.}

    Помните:

    \[
    a^m a^n=a^{m+n},
    \]

    \[
    \frac{a^m}{a^n}=a^{m-n},
    \]

    \[
    (a^m)^n=a^{mn},
    \]

    \[
    a^{-n}=\frac1{a^n}.
    \]

\end{enumerate}

\section*{\textcolor{darkaccent}{11. Короткий чек перед ответом}}

\sectionline

Перед завершением задачи проверьте:

\begin{itemize}[leftmargin=7mm,itemsep=2mm]
    \item правильно ли считаны координаты точек;
    \item все ли неизвестные коэффициенты определены;
    \item верно ли подставлены \(x\) и \(y=f(x)\);
    \item учтено ли условие \(a>0,\ a\ne1\) для показательной функции;
    \item правильно ли преобразованы степени и корни;
    \item собрана ли исходная формула функции;
    \item выполнено ли именно то требование, которое дано в условии.
\end{itemize}

% =========================================================
% ТРЕНИРОВОЧНЫЙ БЛОК
% =========================================================

\clearpage

\section*{\textcolor{darkaccent}{Тренировочный блок --- Путь А}}

\sectionline

\textbf{Количество заданий: 5.}

Решайте задачи методом узловых точек. В тренировочном блоке приведены только условия.

\vspace{3mm}

\subsection*{\textcolor{accent}{Задание 1}}

График функции

\[
f(x)=a^x+b
\]

проходит через точки

\[
A(0;-1),
\qquad
B(1;1).
\]

Найдите

\[
f(4).
\]

\vspace{5mm}

\subsection*{\textcolor{accent}{Задание 2}}

График функции

\[
f(x)=a^x+b
\]

проходит через точки

\[
A(0;-2),
\qquad
B(2;2).
\]

Найдите

\[
f(6).
\]

\vspace{5mm}

\subsection*{\textcolor{accent}{Задание 3}}

Функция имеет вид

\[
f(x)=a^x-4.
\]

Известно, что график проходит через точку

\[
A(-2;0).
\]

Найдите

\[
f(2).
\]

\vspace{5mm}

\subsection*{\textcolor{accent}{Задание 4}}

График функции

\[
f(x)=a^{x+b}
\]

проходит через точки

\[
A(-2;1),
\qquad
B(0;9).
\]

Найдите

\[
f(1).
\]

\vspace{5mm}

\subsection*{\textcolor{accent}{Задание 5}}

График линейной функции

\[
f(x)=kx+b
\]

проходит через точки

\[
A(-1;1),
\qquad
B(2;7).
\]

Найдите значение \(x\), при котором

\[
f(x)=15.
\]

% =========================================================
% ОТВЕТЫ
% =========================================================

\clearpage

\section*{\textcolor{darkaccent}{Ответы}}

\sectionline

\renewcommand{\arraystretch}{1.4}

\begin{table}[ht]
\centering
\begin{tabularx}{0.88\linewidth}{@{}>{\centering\arraybackslash}p{25mm}
>{\centering\arraybackslash}X@{}}
\toprule
\textbf{Задание} & \textbf{Ответ} \\
\midrule
1 & \(79\) \\
2 & \(62\) \\
3 & \(-\dfrac{15}{4}\) \\
4 & \(27\) \\
5 & \(x=6\) \\
\bottomrule
\end{tabularx}
\end{table}

\vfill

\begin{center}
{\small
\textcolor{darkaccent}{
Лёвин Артём Александрович
}}
\end{center}

\end{document}